% =====================================================================
% Determinacy-Field Mechanics, Paper 5 -- DRAFT (v0.1, 2026-08-24)
% v0.1 (第191便): Created by the adopted Paper 4/5 split (原仮定者裁定 2026-08-24
%   「論文4/5分割: 採用」). This paper takes over the MICRO sector of the former
%   Paper 4 draft (dfm-paper4.tex v0.3, "Emergence across scales"):
%   - Sec. "Micro scale: thermalization of the internal-spin temperature" is moved
%     here VERBATIM (prose and measured numbers unchanged; harnesses exp-4-90/91/93).
%   - The emergence standard (E0--E3 grading + six standard tests) is restated in
%     full so this paper stands alone; it was introduced in the former joint draft
%     and is shared verbatim with Paper 4.
%   - The hidden-mass ladder (formerly a Paper-4 macro bullet) moves here per the
%     split line "Paper 4 = gravity/dragging/morphology (macro), Paper 5 =
%     spin/heat/optics (micro)".
%   New sections (core v2 transitions, dimming classes D0--D6, dimming/disappearance
%   design-constraint experiments) are PLANNED skeletons: no measured numbers are
%   claimed there yet, and nothing in them is a claim about real stars.
% Status: draft skeleton with the micro chapter in full prose.
% Companion papers: Paper 1 (foundations), Paper 2 (box universe),
%   Paper 3 (reality calibration), Paper 4 (celestial reproduction).
% =====================================================================
\documentclass[aps,reprint,amsmath,amssymb,nofootinbib,superscriptaddress,floatfix]{revtex4-2}

\usepackage{graphicx}
\usepackage{hyperref}
\usepackage{booktabs}

\begin{document}

\title{Micro-Scale Phase Transitions in Determinacy-Field Mechanics:\\
Internal-Spin Thermodynamics, Core Transitions,\\
and the Optics of Dark Bodies}

% TODO(author): author block as in Papers 1--2.
\author{...}
\date{2026-08-24 (draft v0.1)}

\begin{abstract}
% --- What we claim (skeleton bullets; final prose TODO) ---
% 1. The micro sector of the DFM toy universe -- the internal-spin temperature
%    T_int, its conduction channel, and the spin-driven dimming lS_eff -- is
%    tested with the same graded standard (E0--E3, six standard tests) used
%    across this paper series.
% 2. Micro thermalization (measured, moved from the former Paper-4 draft):
%    inter-group thermalization is a genuine consequence of the conduction
%    channel (frozen exactly when all thermal channels are off), with decay rate
%    lambda ~ lambda_0 (kappa_s/0.15)^{1.006} (n/120)^{1.084}, and an equilibrium
%    temperature predicted by the C*m-weighted ledger.
% 3. Planned: core v2 phase transitions (contraction -> lS_eff saturation,
%    power-ball and cavity states), the dimming ladder D0--D6 as declared optical
%    classes, and design-constraint experiments for deep-dimming and
%    disappearance-like light curves (no identification with any real star).
% 4. The hidden-mass ladder: the model-internal separation of "dark" (optical)
%    from "heavy" (dynamical).
% --- What we do NOT claim ---
%    no real thermodynamics; no identification of any real variable or vanished
%    star; no real dark matter; no 3D extrapolation.
TODO abstract.
\end{abstract}

\maketitle

% =====================================================================
\section{Introduction}

Determinacy-Field Mechanics (DFM) is a two-dimensional relational toy
universe in which a single scalar field -- the determinacy $W$ sourced by
all matter -- supplies the local inertial frame, gravity, clock rates,
and optics \cite{paper1}. It is an educational \emph{computational
thought experiment}: it is not proposed as a description of the real
universe, and none of the phenomena below are claims about real
astrophysical objects (Sec.~\ref{sec:negative}). Paper 1 established the
foundations; Paper 2 confined the model to a finite box universe; Paper 3
calibrated its weak-field sector against real ephemerides under an
explicit ``calibration agreement is not mechanism identification''
philosophy; and Paper 4 pursues celestial reproduction on the macro side
(gravity, dragging, morphology) \cite{paper2,paper3,paper4}.

This paper collects the \emph{micro} side of the same programme: the
thermodynamics of the internal-spin temperature, the phase-transition
behaviour of the core model, and the optics of dark bodies (the dimming
ladder). The split line between Papers 4 and 5, adopted 2026-08-24, is
mechanism-based: Paper 4 = gravity, dragging and morphology (macro);
Paper 5 = spin, heat and optics (micro). The one measured chapter below
(Sec.~\ref{sec:micro}) is moved verbatim from the former joint draft;
the remaining chapters are declared plans whose gates are stated in
advance (Sec.~\ref{sec:standard}).

% =====================================================================
\section{A standard for emergence}
\label{sec:standard}
% [第93便: 本文起こし第1弾 — 方法論章。数値・定義はアプリ docs/PHYSICS.md と
%  ハーネス exp-4-84..93 の実測に一致させてある。査読向け英文推敲は次版]

The word ``emergence'' carries a burden of proof that a single striking
simulation cannot discharge. A fixed random seed, a particular particle
count, and a hand-tuned time window can make almost any dynamical system
produce one beautiful movie. The methodological core of this paper is
therefore a fixed standard, applied identically on every scale we study:
each claimed phenomenon must declare its order parameters in advance,
pass a battery of six standard tests, and carry an explicit grade.

\subsection{Grading}

We grade each phenomenon on a four-level scale, declared per sample and
never auto-assigned:

\begin{itemize}
\item \textbf{E0} -- a configured display: the pattern is placed in the
  initial conditions or enforced by an external element (pinned bodies,
  rails, baths).
\item \textbf{E1} -- a driven effect: the pattern appears while an
  external element does work on the system.
\item \textbf{E2} -- emergence in a closed system: the pattern forms from
  the declared local rules alone, with momentum and angular momentum
  closing on the ledger, but robustness has not been established.
\item \textbf{E3} -- robust emergence: E2, plus reproduction across seeds
  and particle counts and recovery from perturbations, established by the
  standard tests below.
\end{itemize}

The grade is metadata attached to the sample and shown in its user
interface; the measurement harnesses that justify it are cited next to
the grade. Nothing in the engine reads the grade -- it is a claim, not a
mechanism.

\subsection{The six standard tests}

\begin{enumerate}
\item \emph{Knockout control.} Remove one candidate mechanism (a fusion
  rule, a conduction coefficient, a core seed) and re-run under identical
  conditions. If the phenomenon survives, the mechanism was not its
  cause; if every candidate can be removed without effect, the
  phenomenon was configured, not emergent.
\item \emph{Dose response.} Vary the strength of the claimed driver
  monotonically and require the declared order parameter to respond
  monotonically. Non-monotonic responses are reported as such (an
  example: the fusion-threshold sweep of the meso chapter of Paper~4).
\item \emph{Multi-seed.} Re-run under a fixed set of random seeds
  (sixteen in this paper) and report the full range of the order
  parameter, not its best value. Controls run under the same seeds.
\item \emph{Particle-count scaling.} Repeat at several particle counts
  (three to four octaves here) to exclude discrete artefacts of one
  resolution. Where the count is entangled with a physical quantity --
  in our meso-scale sample $n$ is also the total mass -- the entanglement
  is stated rather than hidden.
\item \emph{Perturbation recovery.} Disturb the formed state -- with the
  application's own momentum-conserving injector where possible -- and
  measure whether the order parameters return. A pattern that never
  recovers is a relic of the initial conditions, not an attractor.
\item \emph{Time window.} State the window in which the claim holds and
  what happens outside it. Timescales that grow with system size or
  coupling (both occur below) are reported as part of the result.
\end{enumerate}

Two reporting rules bind all six tests. First, order parameters are
declared before testing; switching the judged quantity after a failed
run is recorded in the harness header when it happens (it happened twice
across this paper pair -- in the micro study below, Sec.~\ref{sec:micro},
and in the meso study of Paper~4 -- and both
switches are documented with the failed first attempt). Second, negative
and partial results are published in the same format as positive ones:
each sample's user interface carries a ``what did not hold'' list, and
this paper reproduces those lists verbatim in
Sec.~\ref{sec:negative}.

Every test in this section runs headless from a version-controlled
harness (Sec.~\ref{sec:repro}); the numbers quoted in
Sec.~\ref{sec:micro} are the outputs of those
harnesses, not of interactive runs.

% =====================================================================
\section{Micro scale: thermalization of the internal-spin temperature}
\label{sec:micro}
% System: two-box gas (cold T_int=0.01 left, hot 6.25 right; 120+120 grains,
%   identical masses/speeds; G=0, kFrame=0). Order parameter: R(t) =
%   DeltaT(t)/DeltaT(0) between the two populations; primary judged quantity is
%   the decay rate lambda (least-squares slope of ln R).
% Measured (tests/exp-4-90.mjs, exp-4-91.mjs; 2026-08-08):
%   - knockout (kappa_s = mu_F = gamma_n = 0): R = 1.0000, lambda = 0 (frozen).
%   - single knockout kappa_s = 0: lambda ~ 0 -- inter-group transport is
%     carried (almost) entirely by conduction; collisions only dissipate
%     within groups.
%   - dose response: kappa_s in {0.05, 0.15, 0.45} -> lambda monotone,
%     nearly proportional.
%   - multi-seed 16: lambda = 2.45e-5 .. 4.20e-5 per step (spread ratio 1.71).
%   - n scaling 60/120/240 per side: lambda monotone increasing.
%   - scaling fit: lambda ~ lambda_0 (kappa_s/0.15)^{1.006} (n/120)^{1.084},
%     lambda_0 = 3.28e-5 per step.
%   - time window: equilibration to R=0.1 extrapolates to ~6.8e4 steps --
%     the conduction timescale is LONG; this is itself a time-window result
%     (the first version of the harness judged t_eq inside a 3000-step window
%     and failed; the switch of order parameter to lambda is recorded in the
%     harness header).
%   - perturbation recovery (5): temperature re-split -> lambda2/lambda1 =
%     0.96/1.15/1.20; velocity kick via the app's official injector ->
%     lambda3/lambda1 = 1.10/1.32/1.32 (all within [0.5, 2]).
%   - equilibrium-temperature prediction (tests/exp-4-93.mjs): landed 2026-08-08;
%     numbers now in prose below.
% Figures: R(t) curves; lambda vs kappa_s and vs n (log-log with slopes) -- TODO.
% [第100便: 本文起こし — 数値は tests/out/exp-4-90/91/93.json の実測そのまま。
%  本サンプルは G=0 の分子系なので第96〜97便の c₀=30 再パラメータ化の影響を受けない
%  (力学 bit 不変) — 旧実測値がそのまま正である]

\subsection{System and order parameter}

The microscopic testbed is deliberately minimal: a closed box of
$120+120$ identical grains, prepared as a cold population on the left
($T_{\mathrm{int}}=0.01$) and a hot one on the right ($T_{\mathrm{int}}=6.25$),
with gravity and frame dragging off ($G=0$, $k_{\mathrm{frame}}=0$).
Temperature is the internal state variable $T_{\mathrm{int}}$ of the A9$''$
revision (heat energy $E=C\,m\,T_{\mathrm{int}}$), so ``thermalization''
here means transport of $T_{\mathrm{int}}$ by the conduction channel
E10$'$ and by collisions -- not equipartition of translational motion.
The declared order parameter is the inter-group contrast
$R(t)=\Delta T(t)/\Delta T(0)$ between the two populations; the judged
quantity is its decay rate $\lambda$, the least-squares slope of
$\ln R(t)$.

That choice of judged quantity is itself a documented switch. The first
version of the harness judged the equilibration time $t_{\mathrm{eq}}$
inside a 3000-step window and \emph{failed in every condition}: the
conduction timescale of this system is long -- extrapolating the
measured decay to $R=0.1$ gives ${\sim}6.8\times10^{4}$ steps. Rather
than widening the window until the test passed, we replaced the judged
quantity by the rate $\lambda$, recorded the failed first attempt in the
harness header (\texttt{tests/exp-4-90.mjs}), and kept the timescale
itself as the time-window result of the battery (test~6).

\subsection{The six tests}

\emph{Knockout (test 1).} With all thermal channels off
($\kappa_s=\mu_F=\gamma_n=0$) the contrast is frozen exactly:
$R=1.0000$, $\lambda=0$ over three seeds. With \emph{only} conduction
removed ($\kappa_s=0$, collisions on) the rate collapses to the noise
scale ($|\lambda| \lesssim 2\times10^{-6}$ per step, sign not even
stable): inter-group transport is carried essentially entirely by the
conduction channel, while collisions merely dissipate within each
group. The phenomenon has exactly one load-bearing mechanism, and
removing it stops the clock.

\emph{Dose response (test 2).} Sweeping $\kappa_s \in \{0.05, 0.15,
0.45\}$ moves $\lambda$ monotonically and nearly proportionally; the
global fit below gives an exponent $1.006$.

\emph{Multi-seed (test 3).} Sixteen seeds give $\lambda = 2.45\times
10^{-5}$--$4.20\times10^{-5}$ per step (spread ratio 1.71, median
$3.28\times10^{-5}$); no seed fails to relax and no seed relaxes
without the channel.

\emph{Particle-count scaling (test 4).} At $n=60/120/240$ per side the
rate increases monotonically. Jointly with the dose sweep the rates
collapse onto
\begin{equation}
\lambda \;\approx\; \lambda_0
\left(\frac{\kappa_s}{0.15}\right)^{1.006}
\left(\frac{n}{120}\right)^{1.084},
\qquad \lambda_0 = 3.28\times10^{-5}\ \text{per step},
\label{eq:microfit}
\end{equation}
i.e.\ proportional to the conduction coefficient and to the grain
density within a few percent (\texttt{tests/exp-4-91.mjs}).

\emph{Perturbation recovery (test 5).} After equilibration has begun we
disturb the state twice: (a) re-splitting the temperatures re-starts a
relaxation with rate ratios $\lambda_2/\lambda_1 = 0.96/1.15/1.20$
across seeds, and (b) a velocity kick through the application's own
momentum-conserving injector gives $\lambda_3/\lambda_1 =
1.10/1.32/1.32$. All ratios sit inside the declared acceptance band
$[0.5, 2]$: the relaxation is not a one-shot property of the prepared
initial state but a rate the mechanism re-establishes after arbitrary
interruptions.

\emph{Time window (test 6).} The honest number is the one that broke
the first harness: full equilibration takes ${\sim}7\times10^{4}$ steps
at the default coupling. We report the rate, and we report that the
timescale is long.

\subsection{A ledger-level prediction}

The battery above establishes robustness; it does not yet show the
sector is \emph{quantitative} inside the model. For that we use the
conduction-only configuration ($\mu_F=\gamma_n=k_{\mathrm{rep}}=0$),
in which the heat $Q=\sum C\,m\,T_{\mathrm{int}}$ should be exactly
conserved and the equilibrium temperature should be the
$C\,m$-weighted mean
\begin{equation}
T_{\mathrm{eq}} \;=\; \frac{\sum_i C_i m_i T_{0,i}}{\sum_i C_i m_i}.
\label{eq:teq}
\end{equation}
Measured (\texttt{tests/exp-4-93.mjs}, two seeds each): with equal
masses the prediction is $T_{\mathrm{eq}}=3.13$ and both populations
approach it monotonically from either side while $Q$ drifts by
${\sim}10^{-6}$ (relative). The discriminating case is unequal masses
(left grains at $m=2$): the naive arithmetic mean would remain $3.13$,
but the ledger predicts $T_{\mathrm{eq}}=2.09$ -- and the measured
populations bracket and approach $2.09$, again with $Q$ conserved to
${\sim}10^{-6}$. The prediction is non-trivial (it distinguishes
mass-weighted from naive averaging) and it is derived from the same
bookkeeping that the conservation monitor displays in the application.
In the default, dissipative configuration the same ledger acquires a
collision-dissipation term; the clean statement is deliberately made in
the configuration where the channel under test is the only one open.

With the six tests and the ledger prediction, the microscopic sector
carries the full battery of Sec.~\ref{sec:standard}. What it
establishes is modest and precise: in this toy universe, inter-group
thermalization is a genuine, robust consequence of one declared channel,
with a rate law [Eq.~(\ref{eq:microfit})] and an equilibrium point
[Eq.~(\ref{eq:teq})] that survive seeds, densities, and perturbations.

% =====================================================================
\section{Core v2 phase transitions (planned)}
\label{sec:corev2}
% PLANNED chapter -- no measured numbers claimed yet.
% Content per the adopted split: the core-v2 state variables (M_c, R_c, J_core)
% under contraction; the saturation of the effective dimming
% lS_eff = min(1, |s| R / c_surf) as the transition marker; the "power-ball"
% and cavity regimes; which order parameters are declared BEFORE measuring.
% Gates: the six standard tests of Sec. II, plus conservation ledgers and
% dt/softening convergence, as in the micro chapter.
TODO.

% =====================================================================
\section{The dimming ladder: declared optical classes D0--D6 (planned)}
\label{sec:dimming}
% PLANNED chapter. D0--D6 are DECLARED classes of the toy universe's dimming
% mechanism (from "bright" to "optically dark while dynamically heavy"),
% not a taxonomy of real objects. The chapter's deliverable is the class
% definitions, their machine-checkable signatures, and negative controls
% (dimming decouples from dynamics bit-exactly -- established in the meso
% chapter of Paper 4).
TODO.

% =====================================================================
\section{Design-constraint experiments: deep dimming and disappearance (planned)}
\label{sec:constraints}
% PLANNED chapter. Two design-constraint experiments ask what parameter regions
% of the declared optical classes could produce (i) deep, irregular dimming
% light curves and (ii) disappearance-like light curves, INSIDE the toy
% universe. These are design constraints on the model's own classes -- explicitly
% NOT identifications of any real variable star or any real vanished source,
% and they are independent of the macro programme of Paper 4 (they gate nothing
% there and are gated by nothing there).
TODO.

% =====================================================================
\section{The hidden-mass ladder}
\label{sec:ladder}
% Moved here from the former Paper-4 macro skeleton per the split line
% (optics/bookkeeping = micro side):
% - Hidden-mass ladder: true mass / luminosity proxy / Newtonian dynamical
%   mass v^2 r / G / frame-drag excess; kFrame=0 control. "Dark" and "heavy"
%   are separated model-internally.
TODO.

% =====================================================================
\section{Limitations and negative claims}
\label{sec:negative}
% Fixed list (kept in spirit from Papers 1--4):
%  1. 2D only; no quantitative extrapolation to 3D.
%  2. No Lorentz invariance; absolute simultaneity; no causal field propagation.
%  3. Dimming is model-specific (not radiative transfer).
%  4. The thermal sector is an analogy; no quantitative correspondence to real
%     thermodynamics (the energy bookkeeping of Sec. III is model-internal).
%  5. The dimming classes D0--D6 are declared classes of the toy universe;
%     no identification with any real variable star or vanished source.
%  6. Dark rotor != dark matter identification (see Paper 4).
%  7. No astrophysical abundance or formation-rate claims.
TODO.

% =====================================================================
\section{Reproducibility}
\label{sec:repro}
% Manifest table: phenomenon -> preset id -> harness -> seed set -> window ->
% commit -> result JSON. Harnesses for the measured chapter: exp-4-90 (micro
% standard tests), exp-4-91 (micro perturbation recovery + scaling fit),
% exp-4-93 (equilibrium-temperature prediction). Planned chapters will declare
% their harnesses before measuring, in the same format.
TODO.

% =====================================================================
\section*{Acknowledgments}
TODO.

% =====================================================================
\begin{thebibliography}{9}
\bibitem{paper1} TODO: Paper 1 reference.
\bibitem{paper2} TODO: Paper 2 reference.
\bibitem{paper3} TODO: Paper 3 reference.
\bibitem{paper4} TODO: Paper 4 reference (celestial reproduction).
\end{thebibliography}

\end{document}
